\bibitem[1]{terui-ref-1} Andrea Cantini. The undecidability of Gris̆in’s set theory. Studia Logica, 74(3):345–368, 2003.
doi:10.1023/A:1025159016268.
\bibitem[2]{terui-ref-2} C. C. Chang. The axiom of comprehension in infinite valued logic. Math. Scand., 13:9–30,
1963. doi:10.7146/math.scand.a-10685.
\bibitem[3]{terui-ref-3} Alonzo Church. A set of postulates for the foundation of logic. Ann. of Math. (2), 33(2):346–366,
1932. doi:10.2307/1968337.
\bibitem[4]{terui-ref-4} Alonzo Church. A set of postulates for the foundation of logic. Ann. of Math. (2), 34(4):839–864,
1933. doi:10.2307/1968702.
\bibitem[5]{terui-ref-5} Alonzo Church. A proof of freedom from contradiction. In Proceedings of the National Academy
of Sciences, volume 21 (5), pages 275–281, 1935. doi:10.1073/pnas.21.5.275.
\bibitem[6]{terui-ref-6} Agata Ciabattoni, Nikolaos Galatos, and Kazushige Terui. Algebraic proof theory for sub-
structural logics: cut-elimination and completions. Ann. Pure Appl. Logic, 163(3):266–290,
2012. doi:10.1016/j.apal.2011.09.003.
\bibitem[7]{terui-ref-7} Agata Ciabattoni, Nikolaos Galatos, and Kazushige Terui. Algebraic proof theory: hy-
persequents and hypercompletions. Ann. Pure Appl. Logic, 168(3):693–737, 2017. doi:
10.1016/j.apal.2016.10.012.
\bibitem[8]{terui-ref-8} Petr Cintula, Petr Hájek, and Carles Noguera, editors. Handbook of mathematical fuzzy
logic. Volume 1, volume 37 of Studies in Logic (London). College Publications, London, 2011.
Mathematical Logic and Foundations.
\bibitem[9]{terui-ref-9} Petr Cintula, Petr Hájek, and Carles Noguera, editors. Handbook of mathematical fuzzy
logic. Volume 2, volume 38 of Studies in Logic (London). College Publications, London, 2011.
Mathematical Logic and Foundations.
\bibitem[10]{terui-ref-10} Abhishek De. A Note on Fixed Points in Łukasiewicz Logic. Log. Univers., 19(4):603–618,
2025. doi:10.1007/s11787-025-00395-1.
\bibitem[11]{terui-ref-11} Jens Erik Fenstad. On the consistency of the axiom of comprehension in the Łukasiewicz
infinite valued logic. Math. Scand., 14:65–74, 1964. doi:10.7146/math.scand.a-10706.
\bibitem[12]{terui-ref-12} Nikolaos Galatos, Peter Jipsen, Tomasz Kowalski, and Hiroakira Ono. Residuated lattices: an
algebraic glimpse at substructural logics, volume 151 of Studies in Logic and the Foundations
of Mathematics. Elsevier B. V., Amsterdam, 2007.
\bibitem[13]{terui-ref-13} Jean-Yves Girard. Light linear logic. Inform. and Comput., 143(2):175–204, 1998. doi:
10.1006/inco.1998.2700.
\bibitem[14]{terui-ref-14} V. N. Grišin. Predicate and set-theoretic calculi based on logic without the contraction rules.
Izv. Akad. Nauk SSSR Ser. Mat., 45(1):47–68, 239, 1981.
\bibitem[15]{terui-ref-15} Petr Hájek. Metamathematics of fuzzy logic, volume 4 of Trends in Logic—Studia Logica
Library. Kluwer Academic Publishers, Dordrecht, 1998. doi:10.1007/978-94-011-5300-3.
\bibitem[16]{terui-ref-16} Petr Hájek. On arithmetic in the Cantor-Łukasiewicz fuzzy set theory. Arch. Math. Logic,
44(6):763–782, 2005. doi:10.1007/s00153-005-0284-0.
\bibitem[17]{terui-ref-17} Petr Hájek. Some remarks on Cantor-Łukasiewicz fuzzy set theory. Log. J. IGPL, 21(2):183–186,
2013. doi:10.1093/jigpal/jzs014.
\bibitem[18]{terui-ref-18} Yuichi Komori. Super-Łukasiewicz propositional logics. Nagoya Math. J., 84:119–133, 1981.
\bibitem[19]{terui-ref-19} Yuichi Komori. Illative combinatory logic based on bck-logic. Mathematica Japonica, 34(4),
1989.
\bibitem[20]{terui-ref-20} Tomasz Kowalski and Tadeusz Litak. Completions of GBL-algebras: negative results. Algebra
Universalis, 58(4):373–384, 2008. doi:10.1007/s00012-008-2056-2.
\bibitem[21]{terui-ref-21} Richard McKinley. Soft linear set theory. J. Log. Algebr. Program., 76(2):226–245, 2008.
doi:10.1016/j.jlap.2008.02.010.
\bibitem[22]{terui-ref-22} Robert McNaughton. A theorem about infinite-valued sentential logic. J. Symbolic Logic,
16:1–13, 1951. doi:10.2307/2268660.
\bibitem[23]{terui-ref-23} George Metcalfe, Nicola Olivetti, and Dov Gabbay. Sequent and hypersequent calculi for
abelian and Łukasiewicz logics. ACM Trans. Comput. Log., 6(3):578–613, 2005. doi:10.1145/
1071596.1071600.
\bibitem[24]{terui-ref-24} George Metcalfe, Nicola Olivetti, and Dov Gabbay. Proof theory for fuzzy logics, volume 36 of
Applied Logic Series. Springer, New York, 2009.
\bibitem[25]{terui-ref-25} Matteo Mio and Alex Simpson. łukasiewicz \ensuremath{\mu}-calculus. Fund. Inform., 150(3-4):317–346, 2017.
doi:10.3233/FI-2017-1472.
\bibitem[26]{terui-ref-26} Paul S. Mostert and Allen L. Shields. On the structure of semigroups on a compact manifold
with boundary. Ann. of Math. (2), 65:117–143, 1957. doi:10.2307/1969668.
\bibitem[27]{terui-ref-27} Hiroakira Ono and Yuichi Komori. Logics without the contraction rule. J. Symbolic Logic,
50(1):169–201, 1985. doi:10.2307/2273798.
\bibitem[28]{terui-ref-28} Uwe Petersen. Logic without contraction as based on inclusion and unrestricted abstraction.
Studia Logica, 64(3):365–403, 2000. doi:10.1023/A:1005293713265.
\bibitem[29]{terui-ref-29} Masaru Shirahata. Linear set theory with strict comprehension. In Proceedings of the Sixth
Asian Logic Conference (Beijing, 1996), pages 223–245. World Sci. Publ., River Edge, NJ,
1998.
\bibitem[30]{terui-ref-30} Thoralf Skolem. Bemerkungen zum Komprehensionsaxiom. Z. Math. Logik Grundlagen Math.,
3:1–17, 1957.
\bibitem[31]{terui-ref-31} Luca Spada. ŁΠ logic with fixed points. Arch. Math. Logic, 47(7-8):741–763, 2008. doi:
10.1007/s00153-008-0105-3.
\bibitem[32]{terui-ref-32} Luca Spada. \ensuremath{\mu}mV-algebras: an approach to fixed points in Łukasiewicz logic. Fuzzy Sets and
Systems, 159(10):1260–1267, 2008. doi:10.1016/j.fss.2007.12.010.
\bibitem[33]{terui-ref-33} Kazushige Terui. Light affine set theory: a naive set theory of polynomial time. Studia Logica,
77(1):9–40, 2004. doi:10.1023/B:STUD.0000034183.33333.6f.
\bibitem[34]{terui-ref-34} Kazushige Terui. A flaw in R.B. White’s article “the consistency of the axiom of comprehension
in the infinite-valued predicate logic of Łukasiewicz”. Unpublished note, 2014. URL: https:
//www.kurims.kyoto-u.ac.jp/~terui/whitenew.pdf.
\bibitem[35]{terui-ref-35} Richard B. White. The consistency of the axiom of comprehension in the infinite-valued
predicate logic of Łukasiewicz. J. Philos. Logic, 8(4):509–534, 1979. doi:10.1007/bf00258447.
\bibitem[36]{terui-ref-36} Richard B. White. A demonstrably consistent type-free extension of the logic BCK. Math.
Japon., 32(1):149–169, 1987.
\bibitem[37]{terui-ref-37} Shunsuke Yatabe. Comprehension contradicts to the induction within Łukasiewicz predicate
logic. Arch. Math. Logic, 48(3-4):265–268, 2009. doi:10.1007/s00153-009-0127-5.
